\nonstopmode{}
\documentclass[a4paper]{book}
\usepackage[times,inconsolata,hyper]{Rd}
\usepackage{makeidx}
\makeatletter\@ifl@t@r\fmtversion{2018/04/01}{}{\usepackage[utf8]{inputenc}}\makeatother
% \usepackage{graphicx} % @USE GRAPHICX@
\makeindex{}
\begin{document}
\chapter*{}
\begin{center}
{\textbf{\huge Package `pkmapr'}}
\par\bigskip{\large \today}
\end{center}
\ifthenelse{\boolean{Rd@use@hyper}}{\hypersetup{pdftitle = {pkmapr: Pakistan Spatial Data Toolkit}}}{}
\ifthenelse{\boolean{Rd@use@hyper}}{\hypersetup{pdfauthor = {Abdullah Umer}}}{}
\begin{description}
\raggedright{}
\item[Title]\AsIs{Pakistan Spatial Data Toolkit}
\item[Version]\AsIs{1.0.0}
\item[Depends]\AsIs{R (>= 4.1.0)}
\item[Description]\AsIs{Provides a tidy interface to Pakistan's official administrative
boundary data from the United Nations Office for the Coordination of
Humanitarian Affairs (OCHA). Downloads and caches spatial data at
country, province, district, and tehsil levels as sf objects
compatible with the tidyverse and geospatial ecosystem. Includes
utilities for geographic dictionary lookup, coordinate reference
system selection, spatial measurement, and neighbour structure
construction for use with spdep, ggplot2, leaflet, and related packages.}
\item[License]\AsIs{GPL-3 + file LICENSE}
\item[Encoding]\AsIs{UTF-8}
\item[Roxygen]\AsIs{list(markdown = TRUE)}
\item[RoxygenNote]\AsIs{7.3.3}
\item[Imports]\AsIs{sf, httr2, dplyr, rlang, cli, spdep, jsonlite}
\item[Suggests]\AsIs{ggplot2, leaflet, rmapshaper, httptest2, testthat (>= 3.0.0),
tmap, knitr, readr, rmarkdown, spelling}
\item[VignetteBuilder]\AsIs{knitr}
\item[URL]\AsIs{}\url{https://abdullahumer1101.github.io/pkmapr}\AsIs{,
}\url{https://github.com/abdullahumer1101/pkmapr}\AsIs{}
\item[BugReports]\AsIs{}\url{https://github.com/abdullahumer1101/pkmapr/issues}\AsIs{}
\item[NeedsCompilation]\AsIs{no}
\item[Author]\AsIs{Abdullah Umer [aut, cre] (ORCID:
<}\url{https://orcid.org/0009-0008-4082-8394}\AsIs{>)}
\item[Maintainer]\AsIs{Abdullah Umer }\email{abdullahumer1101@gmail.com}\AsIs{}
\end{description}
\Rdcontents{Contents}
\HeaderA{get\_country}{Get Pakistan national boundary}{get.Rul.country}
%
\begin{Description}
Get Pakistan national boundary
\end{Description}
%
\begin{Usage}
\begin{verbatim}
get_country(crs = 4326)
\end{verbatim}
\end{Usage}
%
\begin{Arguments}
\begin{ldescription}
\item[\code{crs}] Integer EPSG code. Default 4326 (WGS84). Use 32642 for
distance and area calculations. See \code{\LinkA{pk\_crs\_suggest()}{pk.Rul.crs.Rul.suggest}} for guidance.
\end{ldescription}
\end{Arguments}
%
\begin{Value}
An sf object with columns: country\_name, country\_code,
area\_km2, geometry.
\end{Value}
%
\begin{Examples}
\begin{ExampleCode}

  pk <- get_country()
  plot(sf::st_geometry(pk))

\end{ExampleCode}
\end{Examples}
\HeaderA{get\_districts}{Get Pakistan district boundaries}{get.Rul.districts}
%
\begin{Description}
Get Pakistan district boundaries
\end{Description}
%
\begin{Usage}
\begin{verbatim}
get_districts(province = NULL, crs = 4326)
\end{verbatim}
\end{Usage}
%
\begin{Arguments}
\begin{ldescription}
\item[\code{province}] Character. Filter to one province by exact name.
Matching is case-insensitive. Run \code{pk\_dictionary("provinces")}
to see valid names. NULL (default) returns all districts.

\item[\code{crs}] Integer EPSG code. Default 4326 (WGS84). Use 32642 for
distance and area calculations. See \code{\LinkA{pk\_crs\_suggest()}{pk.Rul.crs.Rul.suggest}} for guidance.
\end{ldescription}
\end{Arguments}
%
\begin{Value}
An sf object with columns: province\_name, district\_name,
district\_code, area\_km2, geometry.
\end{Value}
%
\begin{Examples}
\begin{ExampleCode}

  all_districts    <- get_districts()
  punjab_districts <- get_districts(province = "Punjab")
  punjab_districts <- get_districts(province = "punjab")  # Case-insensitive

\end{ExampleCode}
\end{Examples}
\HeaderA{get\_provinces}{Get Pakistan province boundaries}{get.Rul.provinces}
%
\begin{Description}
Get Pakistan province boundaries
\end{Description}
%
\begin{Usage}
\begin{verbatim}
get_provinces(crs = 4326)
\end{verbatim}
\end{Usage}
%
\begin{Arguments}
\begin{ldescription}
\item[\code{crs}] Integer EPSG code. Default 4326 (WGS84). Use 32642 for
distance and area calculations. See \code{\LinkA{pk\_crs\_suggest()}{pk.Rul.crs.Rul.suggest}} for guidance.
\end{ldescription}
\end{Arguments}
%
\begin{Value}
An sf object with columns: province\_name, province\_code,
area\_km2, geometry.
\end{Value}
%
\begin{Examples}
\begin{ExampleCode}

  pk <- get_provinces()
  plot(sf::st_geometry(pk))

\end{ExampleCode}
\end{Examples}
\HeaderA{get\_tehsils}{Get Pakistan tehsil boundaries}{get.Rul.tehsils}
%
\begin{Description}
Get Pakistan tehsil boundaries
\end{Description}
%
\begin{Usage}
\begin{verbatim}
get_tehsils(district = NULL, province = NULL, crs = 4326)
\end{verbatim}
\end{Usage}
%
\begin{Arguments}
\begin{ldescription}
\item[\code{district}] Character. Filter to one district by exact name.
Matching is case-insensitive. NULL returns all.

\item[\code{province}] Character. Filter to one province by exact name.
Matching is case-insensitive. NULL returns all.
If both district and province are supplied, district takes precedence.

\item[\code{crs}] Integer EPSG code. Default 4326 (WGS84). Use 32642 for
distance and area calculations. See \code{\LinkA{pk\_crs\_suggest()}{pk.Rul.crs.Rul.suggest}} for guidance.
\end{ldescription}
\end{Arguments}
%
\begin{Value}
An sf object with columns: province\_name, district\_name,
tehsil\_name, tehsil\_code, area\_km2, geometry.
\end{Value}
%
\begin{Examples}
\begin{ExampleCode}

  sindh_tehsils  <- get_tehsils(province = "Sindh")
  sindh_tehsils  <- get_tehsils(province = "sindh")  # Case-insensitive
  lahore_tehsils <- get_tehsils(district = "Lahore")
  lahore_tehsils <- get_tehsils(district = "lahore")  # Case-insensitive

\end{ExampleCode}
\end{Examples}
\HeaderA{pkmapr}{pkmapr: Pakistan Spatial Data Toolkit}{pkmapr}
\aliasA{pkmapr-package}{pkmapr}{pkmapr.Rdash.package}
\keyword{internal}{pkmapr}
%
\begin{Description}
Provides a tidy interface to Pakistan's official administrative
boundary data from the United Nations Office for the Coordination of
Humanitarian Affairs (OCHA). Downloads and caches spatial data at
country, province, district, and tehsil levels as \code{sf} objects
compatible with the tidyverse and geospatial ecosystem.
\end{Description}
%
\begin{Section}{Key functions}

\begin{description}

\item[Data loading] 
\code{\LinkA{get\_districts}{get.Rul.districts}}, \code{\LinkA{get\_tehsils}{get.Rul.tehsils}},
\code{\LinkA{get\_provinces}{get.Rul.provinces}}, \code{\LinkA{get\_country}{get.Rul.country}}:
Download administrative boundaries at different levels.

\item[Lookup] 
\code{\LinkA{pk\_dictionary}{pk.Rul.dictionary}}:
Look up official names and PBS geocodes.

\item[Joining] 
\code{\LinkA{pk\_join}{pk.Rul.join}}:
Join external data with administrative units, with match checking.

\item[Spatial weights] 
\code{\LinkA{pk\_neighbors}{pk.Rul.neighbors}}:
Build spatial weights for spdep and spatialreg.

\item[Visualisation] 
\code{\LinkA{pk\_map}{pk.Rul.map}}, \code{\LinkA{pk\_map\_interactive}{pk.Rul.map.Rul.interactive}}:
Quick choropleth maps with ggplot2 or leaflet.


\end{description}

\end{Section}
%
\begin{Section}{Package options}

\code{pkmapr} uses the following global options:
\begin{description}

\item[\code{pkmapr\_use\_cache}] 
Logical. If \code{TRUE}, caches downloaded data in a persistent
user directory. Default \code{FALSE} (uses temporary directory).

\item[\code{pkmapr\_refresh}] 
Logical. Re-download data even if cached copy exists.
Default \code{FALSE}.


\end{description}

\end{Section}
%
\begin{Author}
\strong{Maintainer}: Abdullah Umer \email{abdullahumer1101@gmail.com} (\Rhref{https://orcid.org/0009-0008-4082-8394}{ORCID})

\end{Author}
%
\begin{References}
United Nations Office for the Coordination of Humanitarian Affairs (OCHA).
Pakistan Administrative Boundaries. \url{https://data.humdata.org/}
\end{References}
%
\begin{SeeAlso}
Useful links:
\begin{itemize}

\item{} \Rhref{https://abdullahumer1101.github.io/pkmapr}{Package website}
\item{} \Rhref{https://github.com/abdullahumer1101/pkmapr}{Source code}
\item{} \Rhref{https://github.com/abdullahumer1101/pkmapr/issues}{Bug reports}

\end{itemize}

\end{SeeAlso}
\HeaderA{pkmapr-deprecated}{Deprecated functions in pkmapr}{pkmapr.Rdash.deprecated}
\aliasA{pak\_area}{pkmapr-deprecated}{pak.Rul.area}
\aliasA{pak\_basemap}{pkmapr-deprecated}{pak.Rul.basemap}
\aliasA{pak\_bbox}{pkmapr-deprecated}{pak.Rul.bbox}
\aliasA{pak\_buffer}{pkmapr-deprecated}{pak.Rul.buffer}
\aliasA{pak\_centroid}{pkmapr-deprecated}{pak.Rul.centroid}
\aliasA{pak\_crs\_suggest}{pkmapr-deprecated}{pak.Rul.crs.Rul.suggest}
\aliasA{pak\_dictionary}{pkmapr-deprecated}{pak.Rul.dictionary}
\aliasA{pak\_distance}{pkmapr-deprecated}{pak.Rul.distance}
\aliasA{pak\_intersect}{pkmapr-deprecated}{pak.Rul.intersect}
\aliasA{pak\_join}{pkmapr-deprecated}{pak.Rul.join}
\aliasA{pak\_map}{pkmapr-deprecated}{pak.Rul.map}
\aliasA{pak\_map\_interactive}{pkmapr-deprecated}{pak.Rul.map.Rul.interactive}
\aliasA{pak\_neighbors}{pkmapr-deprecated}{pak.Rul.neighbors}
\aliasA{pak\_points\_in}{pkmapr-deprecated}{pak.Rul.points.Rul.in}
\aliasA{pak\_project}{pkmapr-deprecated}{pak.Rul.project}
\aliasA{pak\_union}{pkmapr-deprecated}{pak.Rul.union}
\keyword{internal}{pkmapr-deprecated}
%
\begin{Description}
These functions are deprecated and will be removed in a future version.
Please use the new \code{pk\_*} functions instead.
\end{Description}
%
\begin{Usage}
\begin{verbatim}
pak_crs_suggest(...)

pak_dictionary(...)

pak_join(...)

pak_area(...)

pak_distance(...)

pak_centroid(...)

pak_bbox(...)

pak_project(...)

pak_neighbors(...)

pak_points_in(...)

pak_buffer(...)

pak_union(...)

pak_intersect(...)

pak_map(...)

pak_map_interactive(...)

pak_basemap(...)
\end{verbatim}
\end{Usage}
\HeaderA{pk\_area}{Recalculate area in km² for an sf object}{pk.Rul.area}
%
\begin{Description}
Computes accurate areas by reprojecting to UTM Zone 42N before
measurement. Returns the input object with area\_km2 added or updated.
\end{Description}
%
\begin{Usage}
\begin{verbatim}
pk_area(x)
\end{verbatim}
\end{Usage}
%
\begin{Arguments}
\begin{ldescription}
\item[\code{x}] An sf object with polygon geometries.
\end{ldescription}
\end{Arguments}
%
\begin{Value}
The input sf object with area\_km2 column added or updated.
\end{Value}
%
\begin{Examples}
\begin{ExampleCode}

  districts <- get_districts()
  districts <- pk_area(districts)

\end{ExampleCode}
\end{Examples}
\HeaderA{pk\_basemap}{Leaflet basemap centred on Pakistan}{pk.Rul.basemap}
%
\begin{Description}
Returns a leaflet map pre-zoomed to Pakistan's extent as a starting
point for building custom interactive maps.
\end{Description}
%
\begin{Usage}
\begin{verbatim}
pk_basemap(provider = "CartoDB.Positron")
\end{verbatim}
\end{Usage}
%
\begin{Arguments}
\begin{ldescription}
\item[\code{provider}] Character. Tile provider. Default "CartoDB.Positron".
\end{ldescription}
\end{Arguments}
%
\begin{Value}
A leaflet object.
\end{Value}
%
\begin{Examples}
\begin{ExampleCode}

  pk_basemap()

\end{ExampleCode}
\end{Examples}
\HeaderA{pk\_bbox}{Get a bounding box for a named administrative unit}{pk.Rul.bbox}
%
\begin{Description}
Get a bounding box for a named administrative unit
\end{Description}
%
\begin{Usage}
\begin{verbatim}
pk_bbox(name, level = c("province", "district", "tehsil"))
\end{verbatim}
\end{Usage}
%
\begin{Arguments}
\begin{ldescription}
\item[\code{name}] Character. Name of the administrative unit.

\item[\code{level}] Character. One of "province", "district", or "tehsil".
\end{ldescription}
\end{Arguments}
%
\begin{Value}
A bbox object for use with \code{ggplot2::coord\_sf()} or
\code{leaflet::fitBounds()}.
\end{Value}
%
\begin{Note}
If you see an error like \AsIs{\texttt{object 'xxx' not found}} when using this
function, the issue is likely in your data preparation, not \code{pk\_bbox()}.
Test the function directly: \code{pk\_bbox("Punjab", level = "province")}.
If that works, check that your data has the expected column names.
\end{Note}
%
\begin{Examples}
\begin{ExampleCode}

  bb <- pk_bbox("Lahore", level = "district")

\end{ExampleCode}
\end{Examples}
\HeaderA{pk\_buffer}{Create buffers around sf geometries in km}{pk.Rul.buffer}
%
\begin{Description}
Reprojects to UTM Zone 42N internally so buffer distances are
in kilometres, then returns buffers in the original CRS.
\end{Description}
%
\begin{Usage}
\begin{verbatim}
pk_buffer(x, dist_km)
\end{verbatim}
\end{Usage}
%
\begin{Arguments}
\begin{ldescription}
\item[\code{x}] An sf object.

\item[\code{dist\_km}] Numeric. Buffer distance in kilometres.
\end{ldescription}
\end{Arguments}
%
\begin{Value}
An sf object with buffered geometries in the same CRS as input.
\end{Value}
%
\begin{Examples}
\begin{ExampleCode}

  districts <- get_districts()
  buffered  <- pk_buffer(districts, dist_km = 10)

\end{ExampleCode}
\end{Examples}
\HeaderA{pk\_centroid}{Extract centroids from an sf object}{pk.Rul.centroid}
%
\begin{Description}
Returns polygon centroids as a point sf object with only geometry
(attributes are dropped to avoid warnings about constant attributes).
\end{Description}
%
\begin{Usage}
\begin{verbatim}
pk_centroid(x)
\end{verbatim}
\end{Usage}
%
\begin{Arguments}
\begin{ldescription}
\item[\code{x}] An sf object with polygon geometries.
\end{ldescription}
\end{Arguments}
%
\begin{Value}
An sf point object in the same CRS as input.
\end{Value}
%
\begin{Examples}
\begin{ExampleCode}

  districts <- get_districts()
  centres   <- pk_centroid(districts)

\end{ExampleCode}
\end{Examples}
\HeaderA{pk\_crs\_suggest}{Suggest an appropriate projected CRS for a Pakistan sf object}{pk.Rul.crs.Rul.suggest}
%
\begin{Description}
Examines the geographic extent of an sf object and recommends the most
appropriate projected coordinate reference system for metric operations.
Pakistan spans UTM zones 41N, 42N, and 43N; national-level analyses
benefit from an equal-area projection.
\end{Description}
%
\begin{Usage}
\begin{verbatim}
pk_crs_suggest(x)
\end{verbatim}
\end{Usage}
%
\begin{Arguments}
\begin{ldescription}
\item[\code{x}] An sf object.
\end{ldescription}
\end{Arguments}
%
\begin{Details}
Using WGS84 (EPSG:4326) for distance, area, or buffer operations produces
inaccurate results as it measures in degrees rather than metres.
\end{Details}
%
\begin{Value}
A named list:
\begin{description}

\item[epsg] Integer EPSG code for the recommended CRS.
\item[name] Human-readable CRS name.
\item[rationale] One-sentence explanation of the recommendation.

\end{description}

\end{Value}
%
\begin{Examples}
\begin{ExampleCode}

  pk_crs_suggest(get_country())
  pk_crs_suggest(get_districts(province = "Balochistan"))

\end{ExampleCode}
\end{Examples}
\HeaderA{pk\_dictionary}{Pakistan Administrative Boundaries Dictionary}{pk.Rul.dictionary}
%
\begin{Description}
Search and explore administrative names and codes for Pakistan.
\end{Description}
%
\begin{Usage}
\begin{verbatim}
pk_dictionary(level = NULL, name = NULL, code = NULL)
\end{verbatim}
\end{Usage}
%
\begin{Arguments}
\begin{ldescription}
\item[\code{level}] Character. One of: "provinces", "districts", "tehsils".
NULL (default) returns all levels.

\item[\code{name}] Character. Filter by partial name (case-insensitive).

\item[\code{code}] Character. Filter by partial P-code (case-insensitive).
\end{ldescription}
\end{Arguments}
%
\begin{Value}
A data frame with columns: name, level, code, parent.
\end{Value}
%
\begin{Section}{Case Insensitivity}

All matching in \code{pk\_dictionary()} is \strong{case-insensitive}.
"Lahore", "lahore", and "LAHORE" all return the same results.
\end{Section}
%
\begin{Examples}
\begin{ExampleCode}
## Not run: 
  # All provinces
  pk_dictionary(level = "provinces")

  # Case-insensitive search for districts containing "lahore"
  pk_dictionary(level = "districts", name = "lahore")
  pk_dictionary(level = "districts", name = "LAHORE")  # Same result

  # Search by code
  pk_dictionary(code = "PK6")

## End(Not run)
\end{ExampleCode}
\end{Examples}
\HeaderA{pk\_distance}{Compute distances between two sf objects in km}{pk.Rul.distance}
%
\begin{Description}
Reprojects internally to UTM Zone 42N for accurate metric distances.
\end{Description}
%
\begin{Usage}
\begin{verbatim}
pk_distance(x, y, by = c("centroid", "edge"))
\end{verbatim}
\end{Usage}
%
\begin{Arguments}
\begin{ldescription}
\item[\code{x}] An sf object.

\item[\code{y}] An sf object.

\item[\code{by}] Character. "centroid" (default) for centroid-to-centroid
distances. "edge" for nearest-point-on-boundary distances.
\end{ldescription}
\end{Arguments}
%
\begin{Value}
A numeric matrix of distances in km.
\end{Value}
%
\begin{Examples}
\begin{ExampleCode}

  provinces <- get_provinces()
  d <- pk_distance(provinces, provinces)

\end{ExampleCode}
\end{Examples}
\HeaderA{pk\_intersect}{Intersect two sf objects}{pk.Rul.intersect}
%
\begin{Description}
Returns the geometric intersection of two sf objects with CRS
alignment handled automatically.
\end{Description}
%
\begin{Usage}
\begin{verbatim}
pk_intersect(x, y)
\end{verbatim}
\end{Usage}
%
\begin{Arguments}
\begin{ldescription}
\item[\code{x}] An sf object.

\item[\code{y}] An sf object.
\end{ldescription}
\end{Arguments}
%
\begin{Value}
An sf object of the intersection in the CRS of x.
\end{Value}
\HeaderA{pk\_join}{Join data to a pkmapr sf object with match checking}{pk.Rul.join}
%
\begin{Description}
Performs a left join. Join by code columns (e.g. district\_code) rather than name columns
wherever possible.
\end{Description}
%
\begin{Usage}
\begin{verbatim}
pk_join(spatial, data, by)
\end{verbatim}
\end{Usage}
%
\begin{Arguments}
\begin{ldescription}
\item[\code{spatial}] An sf object from a pkmapr geometry function.

\item[\code{data}] A data frame to join.

\item[\code{by}] Character. Column name present in both spatial and data.
\end{ldescription}
\end{Arguments}
%
\begin{Value}
The spatial sf object with data columns joined.
\end{Value}
%
\begin{Note}
After joining, always inspect \code{names(result)} to check for column name
conflicts. If your data shares column names with the spatial object (e.g.,
\code{province\_name}, \code{district\_name}), both versions will be preserved with
\code{.x} and \code{.y} suffixes. Rename or select the appropriate columns before
further analysis.
\end{Note}
%
\begin{Examples}
\begin{ExampleCode}

  districts <- get_districts()
  my_data   <- data.frame(district_code = "PK603", value = 42)
  joined    <- pk_join(districts, my_data, by = "district_code")

\end{ExampleCode}
\end{Examples}
\HeaderA{pk\_map}{Quick choropleth map of a pkmapr sf object}{pk.Rul.map}
%
\begin{Description}
Produces a ggplot2 map for rapid exploratory visualisation.
Returns a ggplot object that can be extended with standard ggplot2 layers.
\end{Description}
%
\begin{Usage}
\begin{verbatim}
pk_map(x, fill = NULL, title = NULL, ...)
\end{verbatim}
\end{Usage}
%
\begin{Arguments}
\begin{ldescription}
\item[\code{x}] An sf object.

\item[\code{fill}] Character. Column name to use as fill variable. NULL (default)
produces an outline map.

\item[\code{title}] Character. Map title. NULL for no title.

\item[\code{...}] Additional arguments passed to \code{ggplot2::geom\_sf()}.
\end{ldescription}
\end{Arguments}
%
\begin{Value}
A ggplot object.
\end{Value}
%
\begin{Examples}
\begin{ExampleCode}

  pk_map(get_provinces())
  pk_map(get_provinces(), fill = "area_km2", title = "Province areas")

\end{ExampleCode}
\end{Examples}
\HeaderA{pk\_map\_interactive}{Interactive choropleth map of a pkmapr sf object}{pk.Rul.map.Rul.interactive}
%
\begin{Description}
Produces an interactive leaflet map. Returns a leaflet object that
can be extended with standard leaflet functions.
\end{Description}
%
\begin{Usage}
\begin{verbatim}
pk_map_interactive(x, fill = NULL, popup = NULL, ...)
\end{verbatim}
\end{Usage}
%
\begin{Arguments}
\begin{ldescription}
\item[\code{x}] An sf object.

\item[\code{fill}] Character. Column name for choropleth fill. NULL produces
an outline map.

\item[\code{popup}] Character vector. Column names to display in click popups.

\item[\code{...}] Additional arguments passed to \code{leaflet::addPolygons()}.
\end{ldescription}
\end{Arguments}
%
\begin{Value}
A leaflet object.
\end{Value}
%
\begin{Examples}
\begin{ExampleCode}

  districts <- get_districts()
  pk_map_interactive(districts,
                      fill = "area_km2",
                      popup = list("district_name", "area_km2"))

\end{ExampleCode}
\end{Examples}
\HeaderA{pk\_neighbors}{Construct a spatial neighbours list for Pakistan administrative units}{pk.Rul.neighbors}
%
\begin{Description}
Builds a contiguity or distance-based spatial neighbours structure
for direct use with spdep and spatialreg. Handles Pakistan-specific
complexities including non-contiguous units and disputed boundaries.
\end{Description}
%
\begin{Usage}
\begin{verbatim}
pk_neighbors(
  x,
  style = c("queen", "rook", "knn"),
  k = NULL,
  disputed = c("exclude", "include", "flag")
)
\end{verbatim}
\end{Usage}
%
\begin{Arguments}
\begin{ldescription}
\item[\code{x}] An sf object with polygon geometries.

\item[\code{style}] Character. Neighbour definition: "queen" (shared boundary
point, default), "rook" (shared edge), or "knn" (k nearest centroids).

\item[\code{k}] Integer. Number of nearest neighbours. Required when
\code{style = "knn"}.

\item[\code{disputed}] Character. Treatment of non-contiguous units and
disputed boundaries:
\begin{description}

\item[exclude] Default. Non-contiguous units receive one nearest
neighbour as fallback. An informative message identifies affected units.
\item[include] Treat all boundaries as normal shared boundaries.
\item[flag] Include all boundaries but add a boundary\_note element
to the result documenting which units are affected.

\end{description}

\end{ldescription}
\end{Arguments}
%
\begin{Value}
A named list:
\begin{description}

\item[nb] A spdep nb object.
\item[listw] A row-standardised spdep listw object, ready for
\code{spdep::moran.test()}, \code{spdep::localMoran()},
\code{spatialreg::lagsarlm()}, and related functions.
\item[boundary\_note] Character. Present only when disputed = "flag".

\end{description}

\end{Value}
%
\begin{Section}{Pakistan-specific handling}

Gilgit-Baltistan and Azad Kashmir might break most spatial statistics. The \code{disputed} argument
controls how the Line of Control and special administrative boundaries
flagged in the OCHA source data are treated, making the analytical
decision explicit and reproducible.
\end{Section}
%
\begin{Examples}
\begin{ExampleCode}

  districts <- get_districts()
  w <- pk_neighbors(districts)

  # Calculate Moran's I using spdep
  moran_result <- spdep::moran.test(districts$area_km2, w$listw)

\end{ExampleCode}
\end{Examples}
\HeaderA{pk\_points\_in}{Assign points to administrative units (point-in-polygon)}{pk.Rul.points.Rul.in}
%
\begin{Description}
For each point in \code{points}, identifies which polygon it falls
within and joins that polygon's attributes to the point record.
\end{Description}
%
\begin{Usage}
\begin{verbatim}
pk_points_in(points, polygons, return_all = TRUE)
\end{verbatim}
\end{Usage}
%
\begin{Arguments}
\begin{ldescription}
\item[\code{points}] An sf object with point geometries.

\item[\code{polygons}] An sf object with polygon geometries.

\item[\code{return\_all}] Logical. Keep unmatched points with NA polygon
attributes (TRUE, default) or drop them (FALSE).
\end{ldescription}
\end{Arguments}
%
\begin{Value}
The points sf object with polygon attribute columns joined.
\end{Value}
%
\begin{Examples}
\begin{ExampleCode}

  # Get district boundaries
  districts <- get_districts()

  # Create sample points (or use your own sf object)
  set.seed(123)
  sample_points <- sf::st_sample(districts, size = 50)
  sample_points_sf <- sf::st_sf(geometry = sample_points)

  # Assign points to districts
  points_with_districts <- pk_points_in(sample_points_sf, districts)
  print(head(points_with_districts))

\end{ExampleCode}
\end{Examples}
\HeaderA{pk\_project}{Project an sf object to a Pakistan-appropriate CRS}{pk.Rul.project}
%
\begin{Description}
Convenience wrapper around \code{sf::st\_transform()} with a default
of UTM Zone 42N, appropriate for most Pakistan analyses requiring
accurate distance, area, or buffer operations.
\end{Description}
%
\begin{Usage}
\begin{verbatim}
pk_project(x, crs = 32642)
\end{verbatim}
\end{Usage}
%
\begin{Arguments}
\begin{ldescription}
\item[\code{x}] An sf object.

\item[\code{crs}] Integer EPSG code. Default 32642 (WGS84 / UTM Zone 42N).
\end{ldescription}
\end{Arguments}
%
\begin{Value}
The sf object reprojected to the specified CRS.
\end{Value}
%
\begin{Examples}
\begin{ExampleCode}

  districts <- get_districts()
  projected <- pk_project(districts)

\end{ExampleCode}
\end{Examples}
\HeaderA{pk\_search}{Search Across All Administrative Levels}{pk.Rul.search}
%
\begin{Description}
Search for administrative units by partial name or code across
all levels (provinces, districts, tehsils).
\end{Description}
%
\begin{Usage}
\begin{verbatim}
pk_search(query, fuzzy = FALSE)
\end{verbatim}
\end{Usage}
%
\begin{Arguments}
\begin{ldescription}
\item[\code{query}] Character. Search term (partial match, case-insensitive).

\item[\code{fuzzy}] Logical. If TRUE, uses fuzzy matching for typos.
Default FALSE. Warning: Fuzzy matching can be slower and may
return unexpected matches for short or common queries.
\end{ldescription}
\end{Arguments}
%
\begin{Value}
A data frame with columns: name, level, code, parent.
\end{Value}
%
\begin{Section}{Case Insensitivity}

All matching in \code{pk\_search()} is \strong{case-insensitive} by default.
"Lahore", "lahore", and "LAHORE" all return the same results.
\end{Section}
%
\begin{Section}{Fuzzy Matching}

When \code{fuzzy = TRUE}, the function uses approximate string matching
to handle typos and spelling variations. For example, "lahor" will
match "Lahore". This is useful for interactive exploration but
may return unexpected results for ambiguous queries.
\end{Section}
%
\begin{Examples}
\begin{ExampleCode}
## Not run: 
  # Case-insensitive search
  pk_search("lahore")   # Returns Lahore district and tehsils
  pk_search("LAHORE")   # Same result

  # Fuzzy search for misspelled "lahore"
  pk_search("lahor", fuzzy = TRUE)
  pk_search("lahre", fuzzy = TRUE)

  # Search by code
  pk_search("PK6")

## End(Not run)
\end{ExampleCode}
\end{Examples}
\HeaderA{pk\_union}{Dissolve sf polygons by a grouping column}{pk.Rul.union}
%
\begin{Description}
Merges polygons sharing the same value in a grouping column and
recalculates area\_km2.
\end{Description}
%
\begin{Usage}
\begin{verbatim}
pk_union(x, by)
\end{verbatim}
\end{Usage}
%
\begin{Arguments}
\begin{ldescription}
\item[\code{x}] An sf object.

\item[\code{by}] Character. Column name to group by.
\end{ldescription}
\end{Arguments}
%
\begin{Value}
A dissolved sf object with area\_km2 recalculated.
\end{Value}
%
\begin{Examples}
\begin{ExampleCode}

  tehsils     <- get_tehsils()
  by_district <- pk_union(tehsils, by = "district_name")

\end{ExampleCode}
\end{Examples}
\HeaderA{pk\_version}{Check package version and update status}{pk.Rul.version}
%
\begin{Description}
Check package version and update status
\end{Description}
%
\begin{Usage}
\begin{verbatim}
pk_version(quiet = FALSE)
\end{verbatim}
\end{Usage}
%
\begin{Arguments}
\begin{ldescription}
\item[\code{quiet}] Logical. If FALSE (default), prints to console.
\end{ldescription}
\end{Arguments}
%
\begin{Value}
Invisibly, a list with installed and latest versions.
\end{Value}
%
\begin{Examples}
\begin{ExampleCode}

pk_version()

\end{ExampleCode}
\end{Examples}
\printindex{}
\end{document}
